\chapter{Conclusion}\label{chap:conclusion}

\setlength{\parindent}{1em}

This thesis has presented two systems for automated proof optimization in Lean~4, developed as a natural two-act progression from an agentic prompting approach to a specialized trained model.

\subsection{Summary of Contributions}

\textbf{ImProver} introduced the proof optimization task and demonstrated its feasibility using agentic prompting with a frontier language model. Its key technical contributions are:
\begin{itemize}
    \item The formalization of proof optimization as a task, with a flexible metric framework;
    \item Chain-of-States (CoS) prompting, which annotates intermediate proof states as comments in the proof, providing symbolic information to the LLM;
    \item Error correction and refinement loops that iteratively improve generated proofs;
    \item Retrieval-augmented generation using the Mathlib library and TPiL handbook;
    \item Empirical demonstration that proof optimization generalizes neural theorem proving.
\end{itemize}

\textbf{ImProver~2} addressed the cost and scalability limitations of ImProver by training a small, specialized model. Its key technical contributions are:
\begin{itemize}
    \item An iterative self-improvement pipeline using IRPO with a novel replay buffer that prevents model collapse;
    \item Neurosymbolic augmentation (context extraction, CoS, auto-informalization) that consistently improves performance across model sizes;
    \item Two novel proof optimization metrics: Modularity (maximizing independent subproofs) and Dependencies (minimizing explicitly named external lemmas);
    \item A 7B-parameter model that outperforms models orders-of-magnitude larger and matches frontier models on structural metrics.
\end{itemize}

\subsection{Key Insights}

Several insights emerge from comparing the two systems.

\paragraph{Formal structure exposure is the key ingredient.} Both systems derive their performance from exposing the structure of the Lean proof environment to the model. In ImProver, this is done at inference time via CoS prompting: the model sees intermediate proof states as it generates. In ImProver~2, this is done at training time via neurosymbolic augmentation: the model learns from filtered, augmented preference pairs. The same underlying principle --- giving the model access to formal symbolic data about the proof --- drives performance gains across both paradigms.

\paragraph{Agentic prompting is effective but expensive.} ImProver demonstrates that frontier LLMs, with the right prompting, can perform sophisticated proof optimization. The system achieves strong results on all three datasets and all three metrics, often producing proofs that are substantially shorter or more structured. However, each theorem requires up to 15 LLM calls with document retrieval, making library-scale deployment prohibitively expensive.

\paragraph{Specialization matters more than scale for structural metrics.} ImProver~2 demonstrates that a 7B model, properly trained, outperforms a 671B model on modularity and dependencies. This suggests that structural proof optimization is a specialization problem as much as a scale problem: the bottleneck is not reasoning capacity but whether the model has been adapted to the formal structure of the task.

\paragraph{Reward shaping requires careful guard rails.} Both systems encountered the risk of degenerate solutions: proofs that score highly on a metric without achieving its intent. ImProver addresses this by guiding the model with human-written examples. ImProver~2 uses formal heuristics (duplicate detection, wrapper detection, nontriviality filters) to define the modularity metric in a way that is robust to reward hacking. The design of robust metrics for formal proof quality remains an open challenge.

\subsection{Direct Comparison}

The two systems are complementary rather than competing. On the length metric, ImProver achieves a mean improvement score of $0.355$ on MiniCTX-v2 compared to ImProver~2's $0.330$. This edge is consistent with ImProver's use of an agentic strategy with a frontier model specifically prompted for proof shortening, and its explicit multi-shot example retrieval for each metric. The gap is modest, however, and ImProver's cost per theorem is orders of magnitude higher.

On the structural metrics, ImProver~2 is decisively better. For modularity, ImProver~2 achieves $0.143$ compared to ImProver's $0.088$ --- a $63\%$ improvement. For dependencies, the gap is even larger: $0.206$ for ImProver~2 versus $0.047$ for ImProver. These results demonstrate that iterative structural supervision trains refactoring capabilities that agentic prompting alone does not reliably acquire. A model that has seen thousands of preference-ranked examples of modular rewrites learns what modular structure looks like in a way that in-context examples cannot capture.

The neurosymbolic scaffold (context extraction, CoS, auto-informalization) is beneficial across both systems. Applying the scaffold to a baseline DeepSeek-R1 7B model nearly doubles its length improvement score even before any preference optimization training. Applying it to frontier models yields similarly large gains. This suggests that the scaffold would benefit ImProver's GPT-4o base as well --- a direction not explored in this thesis.

\subsection{Limitations}

\textbf{ImProver} is limited by high cost per theorem (up to 15 GPT-4o calls), dependence on closed-source models, and relatively small evaluation datasets ($72 + 26 + 43$ theorems). These constraints limit both reproducibility and the statistical significance of the results.

\textbf{ImProver~2} is limited by single-step generation (no iterative revision within a single inference call), the fact that the modularity and dependency metrics are structural proxies not yet validated against human maintainer preferences, and training saturation after 2--4 iterations, which limits how much improvement the self-improvement loop can extract from a fixed dataset.

\subsection{Future Directions}

Several natural extensions are motivated by the findings and limitations of this work.

\paragraph{Multi-step revision.} ImProver~2 currently generates improved proofs in a single generation step. Future work could train proof optimizers to perform iterative revision cycles, enabling in-context error fixing similar to ImProver's refinement loop.

\paragraph{Human preference evaluation.} The modularity metric is a structural proxy; whether maintainers actually prefer modularity-optimized proofs remains unknown. A study involving formal mathematicians evaluating the qualitative output of both systems would be valuable for grounding the metrics in real-world utility.

\paragraph{LLM-as-judge metrics.} More subjective or probabilistic metrics --- e.g., using an LLM-as-judge system to score proof readability --- could provide richer optimization targets and better alignment with human preferences.

\paragraph{Combining ImProver and ImProver~2.} ImProver's retrieval strategy provides complementary capabilities to ImProver~2's trained model. A unified system that applies neurosymbolic augmentation at both training and inference time, with additional retrieval, may push performance further.

\paragraph{Application to library maintenance.} Both systems operate on individual theorems; scaling to the full Mathlib library (hundreds of thousands of theorems) and maintaining consistency across inter-dependent rewrites is an open engineering and research challenge.

\paragraph{Fine-tuning downstream provers on optimized data.} A key motivation for proof optimization is improving the quality of training data for downstream neural provers. Empirically validating that provers trained on ImProver- or ImProver~2-optimized corpora outperform those trained on unoptimized data would confirm the training data quality hypothesis.

\subsection{Closing Remarks}

Formal mathematics libraries are growing faster than human maintainers can curate them. This thesis shows that automated proof optimization --- rewriting verified proofs to be shorter, more modular, or less dependent on external lemmas --- is a feasible and learnable task. The central insight is simple but powerful: giving a language model access to the formal structure of the proof environment, whether at inference time through Chain-of-States prompting or at training time through neurosymbolic augmentation, enables far more substantive proof transformations than prompting alone.

ImProver and ImProver~2 together establish proof optimization as a well-defined task with practical applications to library maintenance, training data quality, and the broader agenda of making formal mathematics more accessible to both humans and machines. We hope these systems serve as a foundation for further work on AI-assisted formal mathematics.
